\makeatletter
\def\input@path{{./}{../}{../../}{../../../}{../../../../}{preamble/}{../preamble/}{../../preamble/}{../../../preamble/}}
\makeatother

\input{preamble_exams.tex}
\input{graphs/highlights.tex}
% \GTenHighlightsfalse % Uncomment to disable instructional highlights.
\input{graphs/specific_highlights.tex}
\usepackage{needspace}
\usepackage{titling}
\setlength{\droptitle}{-2em}
\pretitle{\par\begingroup\centering\LARGE}
\posttitle{\par\endgroup}
\preauthor{}\postauthor{}
\predate{}\postdate{}
\maketitlehookd{\vspace{-1em}}
\linespread{1.12}
\setlength{\parskip}{5pt}
\newgeometry{margin=0.65in,top=0.55in,bottom=0.65in}

\title{Grade 11 Term 2 Home Practice -- C1\\\large Parameters and sample statistics\\[0.25em]
	\normalsize Global Economics -- 2026--2027}
\author{}
\date{}

\begin{document}
	\maketitle

	\setcounter{tocdepth}{1}
	\setcounter{secnumdepth}{2}
	\phantomsection
	\label{toc}
	\tableofcontents

	\section*{Evaluated criteria}
	\begin{enumerate}[label=\textbf{C\arabic*},start=1]
		\item Distinguishes population parameters from sample statistics by interpreting and comparing the notation for population and sample means, standard deviations, and proportions in economic and financial contexts.
	\end{enumerate}

	\textbf{How to practise.} Try each task before opening its blue solution link. The arrows return to the contents, problem list or current problem.

	\section{Introduction}
	\label{sec:c1-introduction}

	\SectionProblemsTableOfContents{
		\SectionProblemLink{subsec:c1-population-and-sample}{From electricity bills to a mean} \\
		\SectionProblemLink{subsec:c1-quick-check}{Quick check: redraw or redefine?}
	}

	\subsection{From electricity bills to a mean}
	\label{subsec:c1-population-and-sample}

	A utility wants to describe May electricity bills for small shops in one district. The \textbf{target population} is \HighlightStroke{G11T2C1Population}{all May bills for small shops in that district}. The utility reviews \HighlightStroke{G11T2C1Sample}{40 randomly selected bills} from that list. Those 40 bills form the \textbf{sample}, and the amount charged per bill is the measured variable.

	The mean amount charged across the full list, $\HighlightStroke{G11T2C1PopulationMean}{\mu}$, is a \textbf{population parameter}. The mean across the 40 selected bills, $\HighlightStroke{G11T2C1SampleMean}{\bar{x}}=\HighlightStroke{G11T2C1SampleMean}{210}$ thousand COP, is a \textbf{sample statistic}. The sample mean estimates the population mean. The value of $\mu$ is unknown from the information given.

	\secInicio[subsec:c1-population-and-sample][sec:c1-introduction]

	\subsection{Quick check: redraw or redefine?}
	\label{subsec:c1-quick-check}

	\QuestionsTasks[sec:c1-introduction]{
		{1}/{sol:c1-quick-check}/
		{The utility could draw 40 bills again from the same May list or extend the study to May bills for small shops in two districts and draw 40 from the combined list. Which mean is being estimated in each case?}
	}

	\subsubsection*{Suggested answer}
	\label{sol:c1-quick-check}

	A redraw from the original list can give a different $\HighlightStroke{G11T2C1SampleMean}{\bar{x}}$, while $\HighlightStroke{G11T2C1PopulationMean}{\mu}$ still describes all May bills for small shops in the original district. With the combined list, the new sample estimates the mean for May bills for small shops across both districts. The definition of $\mu$ changes, even though both samples have 40 bills. The two population means could have the same numerical value.

	\secInicio[subsec:c1-quick-check][sec:c1-introduction]

	\clearpage
	\section{Notation Reference and practice procedure}
	\label{sec:c1-notation-and-procedure}

	\SectionProblemsTableOfContents{
		\SectionProblemLink{subsec:c1-notation}{Notation reference} \\
		\SectionProblemLink{subsec:c1-procedure}{Practice procedure}
	}

	\subsection{Notation reference}
	\label{subsec:c1-notation}

	\begin{center}
	\begin{tabularx}{\linewidth}{@{}*{2}{>{\raggedright\arraybackslash}X}>{\raggedright\arraybackslash}p{.23\linewidth}@{}}
	\toprule
	\textbf{Population parameter} & \textbf{Sample statistic} & \textbf{Units} \\
	\midrule
	$\HighlightStroke{G11T2C1PopulationMean}{\mu}$, mean of the variable across all units & $\HighlightStroke{G11T2C1SampleMean}{\bar{x}}$, mean across the selected units & Same as the variable \\
	$\HighlightStroke{G11T2C1PopulationSD}{\sigma}$, standard deviation of the variable across all units & $\HighlightStroke{G11T2C1SampleSD}{s}$, standard deviation across the selected units & Same as the variable \\
	$\HighlightStroke{G11T2C1PopulationProportion}{p}$, fraction of all units meeting a condition & $\HighlightStroke{G11T2C1SampleProportion}{\hat{p}}$, fraction of selected units meeting that condition & Unitless, as a fraction, decimal or percentage \\
	\bottomrule
	\end{tabularx}
	\end{center}

	The mean describes the average level, and the standard deviation describes how spread out the values are. In the utility example, $\mu$ and $\sigma$ describe the amount charged, in thousand COP, across the specified May bills. A different month, district or group of customers changes which bills define the parameters. Measuring electricity use on the same bills would instead give a mean and standard deviation in kilowatt-hours. For a proportion, specify the condition as well as the population. For instance, $p$ could be the fraction of those May bills paid by the due date.

	\secInicio[subsec:c1-notation][sec:c1-notation-and-procedure]

	\subsection{Practice procedure}
	\label{subsec:c1-procedure}

	\begin{enumerate}
		\item Name the unit being studied and define the \HighlightStroke{G11T2C1Population}{target population}, including any time, place or eligibility limits. Then identify the \HighlightStroke{G11T2C1Sample}{selected units} and the sample size.
		\item State the variable and its units for a mean or standard deviation. For a proportion, state the condition and identify the denominator as the size of the group being described. To find $\hat{p}$, divide the number of selected units meeting the condition by the sample size.
		\item Match the quantity to its symbol. Ask whether it summarizes the whole target population or the selected units. Write the contextual meaning even when the numerical value is unknown.
		\item Before subtracting, check that the variable or condition and the units match. Name the two groups and say whether the result compares two sample statistics, a sample statistic with a population parameter, or two known population parameters. Express differences between percentages in percentage points.
		\item For a redraw from the same population with the same measures, the parameters stay fixed and the statistics can vary. If the target population changes, redefine the parameters for the new group. A difference between sample statistics alone cannot establish the size or direction of a difference between population parameters.
	\end{enumerate}

	\secInicio[subsec:c1-procedure][sec:c1-notation-and-procedure]

	\clearpage
	\section{General Worked problems}
	\label{sec:worked-solution-examples}

	\SectionProblemsTableOfContents{
		\SectionProblemLink{prob:c1-1}{A market association's sales survey (C1)} \\
		\SectionProblemLink{prob:c1-2}{A supermarket's checkout survey (C1)} \\
		\SectionProblemLink{prob:c1-3}{A transit agency's fare review (C1)} \\
		\SectionProblemLink{prob:c1-4}{A coffee association's export shipments (C1)} \\
		\SectionProblemLink{prob:c1-5}{A credit union's loan review (C1)}
	}

	% ----------------------------------------------------------------
	\subsection{A market association's sales survey (C1)}
	\label{prob:c1-1}

	A market association wants to describe all 1,200 stalls licensed for June. From 60 randomly selected stalls, it finds a mean June sales revenue of 4.8 million COP and a sample standard deviation of 1.2 million COP. Eighteen sampled stalls accept digital payment.

	\QuestionsTasks[sec:worked-solution-examples]{
		{1}/{sol:c1-1-a}/
		{Identify the population and sample. Give the symbols and contextual meanings for the population mean, standard deviation and proportion, and pair them with the corresponding sample statistics.},
		{1}/{sol:c1-1-b}/
		{To check how the survey results vary across samples, the association draws another random sample from the same June list. It finds a mean of 5.1 million COP, a standard deviation of 1.0 million COP and a digital-payment share of 0.25. Revisit your Task 1 answers. What changes, what stays the same, and which numerical differences can you calculate?}
	}

	\clearpage

	\PracticeHighlightLegend{G11-T2-C1}

	\Needspace{5\baselineskip}
	\phantomsection
	\label{sol:c1-1-a}
	\subsubsection*{C1 -- Match each quantity to its symbol}

	The population is \HighlightStroke{G11T2C1Population}{all 1,200 June-licensed stalls}. The sample is \HighlightStroke{G11T2C1Sample}{the 60 selected stalls}. Each pair describes the same feature at a different scale.

	\begin{center}
	\begin{tabularx}{\linewidth}{@{}*{2}{>{\raggedright\arraybackslash}X}@{}}
	\toprule
	Population parameter & Sample statistic \\
	\midrule
	$\HighlightStroke{G11T2C1PopulationMean}{\mu}$, mean June sales across all licensed stalls & $\HighlightStroke{G11T2C1SampleMean}{\bar{x}}=\HighlightStroke{G11T2C1SampleMean}{4.8}$ million COP, mean across the 60 \\
	$\HighlightStroke{G11T2C1PopulationSD}{\sigma}$, SD of June sales across all licensed stalls & $\HighlightStroke{G11T2C1SampleSD}{s}=\HighlightStroke{G11T2C1SampleSD}{1.2}$ million COP, SD across the 60 \\
	$\HighlightStroke{G11T2C1PopulationProportion}{p}$, fraction of licensed stalls accepting digital payment & $\HighlightStroke{G11T2C1SampleProportion}{\hat p}=18/60=\HighlightStroke{G11T2C1SampleProportion}{0.30}$, fraction among the 60 \\
	\bottomrule
	\end{tabularx}
	\end{center}

	The two sales measures have units of million COP. The payment proportions are unitless and can also be written as percentages.

	\secInicio[prob:c1-1][sec:worked-solution-examples]

	\Needspace{5\baselineskip}
	\phantomsection
	\label{sol:c1-1-b}
	\subsubsection*{C1 -- Compare what changes across samples}


	\begin{center}
	\begin{tabularx}{\linewidth}{@{}>{\raggedright\arraybackslash}p{.16\linewidth}*{2}{>{\raggedright\arraybackslash}p{.22\linewidth}}>{\raggedright\arraybackslash}X@{}}
	\toprule
	Quantity & Original scenario & Revised scenario & What the change tells us \\
	\midrule
	Population & \HighlightStroke{G11T2C1Population}{All 1,200} June-licensed stalls & \HighlightStroke{G11T2C1Population}{Same June list} & $\HighlightStroke{G11T2C1PopulationMean}{\mu}$, $\HighlightStroke{G11T2C1PopulationSD}{\sigma}$ and $\HighlightStroke{G11T2C1PopulationProportion}{p}$ stay fixed. \\
	Sample & \HighlightStroke{G11T2C1Sample}{First 60} selected stalls & Second \HighlightStroke{G11T2C1Sample}{random sample} from the June list & A new draw can produce different statistics. \\
	Mean & $\HighlightStroke{G11T2C1SampleMean}{\bar{x}}=\HighlightStroke{G11T2C1SampleMean}{4.8}$ million COP & $\HighlightStroke{G11T2C1SampleMean}{\bar{x}}=\HighlightStroke{G11T2C1SampleMean}{5.1}$ million COP & $5.1-4.8=+0.3$ million COP in $\HighlightStroke{G11T2C1SampleMean}{\bar{x}}$. \\
	Standard\newline deviation & $\HighlightStroke{G11T2C1SampleSD}{s}=\HighlightStroke{G11T2C1SampleSD}{1.2}$ million COP & $\HighlightStroke{G11T2C1SampleSD}{s}=\HighlightStroke{G11T2C1SampleSD}{1.0}$ million COP & $1.0-1.2=-0.2$ million COP in $\HighlightStroke{G11T2C1SampleSD}{s}$. \\
	Digital-payment share & $\HighlightStroke{G11T2C1SampleProportion}{\hat p}=18/60=\HighlightStroke{G11T2C1SampleProportion}{0.30}$ & $\HighlightStroke{G11T2C1SampleProportion}{\hat p}=\HighlightStroke{G11T2C1SampleProportion}{0.25}$ & $0.25-0.30=-0.05$ in $\HighlightStroke{G11T2C1SampleProportion}{\hat p}$ ($-5$ percentage points). \\
	\bottomrule
	\end{tabularx}
	\end{center}

	\secInicio[prob:c1-1][sec:worked-solution-examples]

	\clearpage

	% ----------------------------------------------------------------
	\subsection{A supermarket's checkout survey (C1)}
	\label{prob:c1-2}

	A supermarket manager wants to describe all 1,500 purchases made at one store on Saturday. In a random sample of 50 receipts, the mean purchase value is 64 thousand COP and the sample standard deviation is 18 thousand COP. Twenty of those purchases used a loyalty card.

	\QuestionsTasks[sec:worked-solution-examples]{
		{1}/{sol:c1-2-a}/
		{Identify the population and sample. Give the symbols and contextual meanings for the population mean, standard deviation and proportion, and pair them with the corresponding sample statistics.},
		{1}/{sol:c1-2-b}/
		{Before reporting the results, the manager asks a colleague to draw another random sample of 50 receipts from that store's Saturday records. The new sample has a mean purchase value of 61 thousand COP and a standard deviation of 20 thousand COP. Twenty-five of these purchases used a loyalty card. Revisit your Task 1 answers. What changes, what stays the same, and which numerical differences can you calculate?}
	}

	\clearpage

	\Needspace{5\baselineskip}
	\phantomsection
	\label{sol:c1-2-a}
	\subsubsection*{C1 -- Match each quantity to its symbol}
	The population is \HighlightStroke{G11T2C1Population}{all 1,500 Saturday purchases} at the store. The sample is \HighlightStroke{G11T2C1Sample}{the 50 purchases} represented by the selected receipts.

	\begin{center}
	\begin{tabularx}{\linewidth}{@{}*{2}{>{\raggedright\arraybackslash}X}@{}}
	\toprule
	Population parameter & Sample statistic \\
	\midrule
	$\HighlightStroke{G11T2C1PopulationMean}{\mu}$, mean value of all Saturday purchases & $\HighlightStroke{G11T2C1SampleMean}{\bar{x}}=\HighlightStroke{G11T2C1SampleMean}{64}$ thousand COP, mean among the 50 purchases \\
	$\HighlightStroke{G11T2C1PopulationSD}{\sigma}$, standard deviation of value across all Saturday purchases & $\HighlightStroke{G11T2C1SampleSD}{s}=\HighlightStroke{G11T2C1SampleSD}{18}$ thousand COP, standard deviation among the 50 purchases \\
	$\HighlightStroke{G11T2C1PopulationProportion}{p}$, fraction of all Saturday purchases that used a loyalty card & $\HighlightStroke{G11T2C1SampleProportion}{\hat p}=20/50=\HighlightStroke{G11T2C1SampleProportion}{0.40}$, fraction among the 50 purchases \\
	\bottomrule
	\end{tabularx}
	\end{center}

	\secInicio[prob:c1-2][sec:worked-solution-examples]

	\Needspace{18\baselineskip}
	\phantomsection
	\label{sol:c1-2-b}
	\subsubsection*{C1 -- Compare what changes across samples}

	\begin{center}
	\begin{tabularx}{\linewidth}{@{}>{\raggedright\arraybackslash}p{.16\linewidth}*{2}{>{\raggedright\arraybackslash}p{.22\linewidth}}>{\raggedright\arraybackslash}X@{}}
	\toprule
	Quantity & Original scenario & Revised scenario & What the change tells us \\
	\midrule
	Population & \HighlightStroke{G11T2C1Population}{All 1,500} Saturday purchases at the store & \HighlightStroke{G11T2C1Population}{Same 1,500 purchases} & $\HighlightStroke{G11T2C1PopulationMean}{\mu}$, $\HighlightStroke{G11T2C1PopulationSD}{\sigma}$ and $\HighlightStroke{G11T2C1PopulationProportion}{p}$ stay fixed. \\
	Sample & \HighlightStroke{G11T2C1Sample}{First draw} of 50 receipts & \HighlightStroke{G11T2C1Sample}{Second draw} of 50 receipts & The sample size stays the same. The draw is new. \\
	Mean & $\HighlightStroke{G11T2C1SampleMean}{\bar{x}}=\HighlightStroke{G11T2C1SampleMean}{64}$ thousand COP & $\HighlightStroke{G11T2C1SampleMean}{\bar{x}}=\HighlightStroke{G11T2C1SampleMean}{61}$ thousand COP & $61-64=-3$ thousand COP in $\HighlightStroke{G11T2C1SampleMean}{\bar{x}}$. \\
	Standard\newline deviation & $\HighlightStroke{G11T2C1SampleSD}{s}=\HighlightStroke{G11T2C1SampleSD}{18}$ thousand COP & $\HighlightStroke{G11T2C1SampleSD}{s}=\HighlightStroke{G11T2C1SampleSD}{20}$ thousand COP & $20-18=+2$ thousand COP in $\HighlightStroke{G11T2C1SampleSD}{s}$. \\
	Loyalty-card share & $\HighlightStroke{G11T2C1SampleProportion}{\hat p}=20/50=\HighlightStroke{G11T2C1SampleProportion}{0.40}$ & $\HighlightStroke{G11T2C1SampleProportion}{\hat p}=25/50=\HighlightStroke{G11T2C1SampleProportion}{0.50}$ & $0.50-0.40=+0.10$ in $\HighlightStroke{G11T2C1SampleProportion}{\hat p}$ ($+10$ percentage points). \\
	\bottomrule
	\end{tabularx}
	\end{center}

	\secInicio[prob:c1-2][sec:worked-solution-examples]

	\clearpage

	% ----------------------------------------------------------------
	\subsection{A transit agency's fare review (C1)}
	\label{prob:c1-3}

	A transit agency wants to describe all 8,400 bus boardings on its central routes last Monday. In a random sample of 100 boarding records, the mean fare paid is 3.2 thousand COP and the sample standard deviation is 0.8 thousand COP. Thirty-five of the sampled boardings used a discounted fare.

	\QuestionsTasks[sec:worked-solution-examples]{
		{1}/{sol:c1-3-a}/
		{Identify the population and sample. Give the symbols and contextual meanings for the population mean, standard deviation and proportion, and pair them with the corresponding sample statistics.},
		{1}/{sol:c1-3-b}/
		{For a fare review, the agency's director asks to include suburban routes alongside the central routes, using the same Monday records. The combined routes have 12,000 boardings. A random sample of 100 from the combined records gives a mean fare of 3.4 thousand COP and a standard deviation of 0.9 thousand COP. Thirty sampled boardings used a discounted fare. Revisit your Task 1 answers. What changes, what stays the same, and which numerical differences can you calculate?}
	}

	\clearpage

	\Needspace{5\baselineskip}
	\phantomsection
	\label{sol:c1-3-a}
	\subsubsection*{C1 -- Match each quantity to its symbol}
	The population is \HighlightStroke{G11T2C1Population}{all 8,400 Monday boardings} on the central routes. The sample is \HighlightStroke{G11T2C1Sample}{the 100 boardings} represented by the selected records.

	\begin{center}
	\begin{tabularx}{\linewidth}{@{}*{2}{>{\raggedright\arraybackslash}X}@{}}
	\toprule
	Population parameter & Sample statistic \\
	\midrule
	$\HighlightStroke{G11T2C1PopulationMean}{\mu}$, mean fare paid across all central-route boardings that Monday & $\HighlightStroke{G11T2C1SampleMean}{\bar{x}}=\HighlightStroke{G11T2C1SampleMean}{3.2}$ thousand COP, mean among the 100 boardings \\
	$\HighlightStroke{G11T2C1PopulationSD}{\sigma}$, standard deviation of fares across all those boardings & $\HighlightStroke{G11T2C1SampleSD}{s}=\HighlightStroke{G11T2C1SampleSD}{0.8}$ thousand COP, standard deviation among the 100 boardings \\
	$\HighlightStroke{G11T2C1PopulationProportion}{p}$, fraction of all those boardings that used a discounted fare & $\HighlightStroke{G11T2C1SampleProportion}{\hat p}=35/100=\HighlightStroke{G11T2C1SampleProportion}{0.35}$, fraction among the 100 boardings \\
	\bottomrule
	\end{tabularx}
	\end{center}

	\secInicio[prob:c1-3][sec:worked-solution-examples]

	\Needspace{18\baselineskip}
	\phantomsection
	\label{sol:c1-3-b}
	\subsubsection*{C1 -- Compare what changes when routes are added}

	\begin{center}
	\begin{tabularx}{\linewidth}{@{}>{\raggedright\arraybackslash}p{.16\linewidth}*{2}{>{\raggedright\arraybackslash}p{.22\linewidth}}>{\raggedright\arraybackslash}X@{}}
	\toprule
	Quantity & Original scenario & Revised scenario & What the change tells us \\
	\midrule
	Population & \HighlightStroke{G11T2C1Population}{8,400} Monday boardings on central routes & \HighlightStroke{G11T2C1Population}{12,000} Monday boardings on central and suburban routes & The population gains $12{,}000-8{,}400=3{,}600$ boardings. The day stays the same. \\
	Sample & \HighlightStroke{G11T2C1Sample}{100 central-route} boardings & \HighlightStroke{G11T2C1Sample}{100 boardings} from the combined routes & The sample size stays at 100. The pool it is drawn from expands. \\
	Mean & $\HighlightStroke{G11T2C1PopulationMean}{\mu}$ for central routes, $\HighlightStroke{G11T2C1SampleMean}{\bar{x}}=\HighlightStroke{G11T2C1SampleMean}{3.2}$ thousand COP & $\HighlightStroke{G11T2C1PopulationMean}{\mu}$ for combined routes, $\HighlightStroke{G11T2C1SampleMean}{\bar{x}}=\HighlightStroke{G11T2C1SampleMean}{3.4}$ thousand COP & $3.4-3.2=+0.2$ thousand COP in $\HighlightStroke{G11T2C1SampleMean}{\bar{x}}$. Whether $\HighlightStroke{G11T2C1PopulationMean}{\mu}$ changed is unknown. \\
	Standard\newline deviation & $\HighlightStroke{G11T2C1PopulationSD}{\sigma}$ for central routes, $\HighlightStroke{G11T2C1SampleSD}{s}=\HighlightStroke{G11T2C1SampleSD}{0.8}$ thousand COP & $\HighlightStroke{G11T2C1PopulationSD}{\sigma}$ for combined routes, $\HighlightStroke{G11T2C1SampleSD}{s}=\HighlightStroke{G11T2C1SampleSD}{0.9}$ thousand COP & $0.9-0.8=+0.1$ thousand COP in $\HighlightStroke{G11T2C1SampleSD}{s}$. Whether $\HighlightStroke{G11T2C1PopulationSD}{\sigma}$ changed is unknown. \\
	Discounted-fare share & $\HighlightStroke{G11T2C1PopulationProportion}{p}$ for central routes, $\HighlightStroke{G11T2C1SampleProportion}{\hat p}=35/100=\HighlightStroke{G11T2C1SampleProportion}{0.35}$ & $\HighlightStroke{G11T2C1PopulationProportion}{p}$ for combined routes, $\HighlightStroke{G11T2C1SampleProportion}{\hat p}=30/100=\HighlightStroke{G11T2C1SampleProportion}{0.30}$ & $0.30-0.35=-0.05$ in $\HighlightStroke{G11T2C1SampleProportion}{\hat p}$ ($-5$ percentage points). Whether $\HighlightStroke{G11T2C1PopulationProportion}{p}$ changed is unknown. \\
	\bottomrule
	\end{tabularx}
	\end{center}

	\secInicio[prob:c1-3][sec:worked-solution-examples]

	\clearpage

	% ----------------------------------------------------------------
	\subsection{A coffee association's export shipments (C1)}
	\label{prob:c1-4}

	A coffee exporters' association wants to describe all 900 coffee shipments in its April export records. In a random sample of 40 shipments, the mean shipment weight is 18.4 metric tonnes and the sample standard deviation is 4.0 metric tonnes. Twelve of the sampled shipments departed on schedule.

	\QuestionsTasks[sec:worked-solution-examples]{
		{1}/{sol:c1-4-a}/
		{Identify the population and sample. Give the symbols and contextual meanings for the population mean, standard deviation and proportion, and pair them with the corresponding sample statistics.},
		{1}/{sol:c1-4-b}/
		{A second analyst repeats the study with a fresh random draw of 40 shipments from the association's April records. The mean shipment weight is 19.1 metric tonnes and the standard deviation is 3.6 metric tonnes. Twelve of these shipments departed on schedule. Revisit your Task 1 answers. What changes, what stays the same, and which numerical differences can you calculate?}
	}

	\clearpage

	\Needspace{5\baselineskip}
	\phantomsection
	\label{sol:c1-4-a}
	\subsubsection*{C1 -- Match each quantity to its symbol}
	The population is \HighlightStroke{G11T2C1Population}{all 900 coffee shipments} in the association's April export records. The sample is \HighlightStroke{G11T2C1Sample}{the 40} selected shipments.

	\begin{center}
	\begin{tabularx}{\linewidth}{@{}*{2}{>{\raggedright\arraybackslash}X}@{}}
	\toprule
	Population parameter & Sample statistic \\
	\midrule
	$\HighlightStroke{G11T2C1PopulationMean}{\mu}$, mean weight across all 900 April shipments & $\HighlightStroke{G11T2C1SampleMean}{\bar{x}}=\HighlightStroke{G11T2C1SampleMean}{18.4}$ metric tonnes, mean among the 40 shipments \\
	$\HighlightStroke{G11T2C1PopulationSD}{\sigma}$, standard deviation of weight across all 900 April shipments & $\HighlightStroke{G11T2C1SampleSD}{s}=\HighlightStroke{G11T2C1SampleSD}{4.0}$ metric tonnes, standard deviation among the 40 shipments \\
	$\HighlightStroke{G11T2C1PopulationProportion}{p}$, fraction of all 900 April shipments that departed on schedule & $\HighlightStroke{G11T2C1SampleProportion}{\hat p}=12/40=\HighlightStroke{G11T2C1SampleProportion}{0.30}$, fraction among the 40 shipments \\
	\bottomrule
	\end{tabularx}
	\end{center}

	\secInicio[prob:c1-4][sec:worked-solution-examples]

	\Needspace{18\baselineskip}
	\phantomsection
	\label{sol:c1-4-b}
	\subsubsection*{C1 -- Compare what changes across samples}

	\begin{center}
	\begin{tabularx}{\linewidth}{@{}>{\raggedright\arraybackslash}p{.16\linewidth}*{2}{>{\raggedright\arraybackslash}p{.22\linewidth}}>{\raggedright\arraybackslash}X@{}}
	\toprule
	Quantity & Original scenario & Revised scenario & What the change tells us \\
	\midrule
	Population & \HighlightStroke{G11T2C1Population}{All 900 shipments} in the April records & \HighlightStroke{G11T2C1Population}{Same 900 shipments} & $\HighlightStroke{G11T2C1PopulationMean}{\mu}$, $\HighlightStroke{G11T2C1PopulationSD}{\sigma}$ and $\HighlightStroke{G11T2C1PopulationProportion}{p}$ stay fixed. \\
	Sample & \HighlightStroke{G11T2C1Sample}{First draw} of 40 shipments & \HighlightStroke{G11T2C1Sample}{Second draw} of 40 shipments & The sample size stays the same. The draw is new. \\
	Mean & $\HighlightStroke{G11T2C1SampleMean}{\bar{x}}=\HighlightStroke{G11T2C1SampleMean}{18.4}$ metric tonnes & $\HighlightStroke{G11T2C1SampleMean}{\bar{x}}=\HighlightStroke{G11T2C1SampleMean}{19.1}$ metric tonnes & $19.1-18.4=+0.7$ metric tonnes in $\HighlightStroke{G11T2C1SampleMean}{\bar{x}}$. \\
	Standard\newline deviation & $\HighlightStroke{G11T2C1SampleSD}{s}=\HighlightStroke{G11T2C1SampleSD}{4.0}$ metric tonnes & $\HighlightStroke{G11T2C1SampleSD}{s}=\HighlightStroke{G11T2C1SampleSD}{3.6}$ metric tonnes & $3.6-4.0=-0.4$ metric tonnes in $\HighlightStroke{G11T2C1SampleSD}{s}$. \\
	On-schedule share & $\HighlightStroke{G11T2C1SampleProportion}{\hat p}=12/40=\HighlightStroke{G11T2C1SampleProportion}{0.30}$ & $\HighlightStroke{G11T2C1SampleProportion}{\hat p}=12/40=\HighlightStroke{G11T2C1SampleProportion}{0.30}$ & $0.30-0.30=0$ in $\HighlightStroke{G11T2C1SampleProportion}{\hat p}$ (0 percentage points). \\
	\bottomrule
	\end{tabularx}
	\end{center}

	\secInicio[prob:c1-4][sec:worked-solution-examples]

	\clearpage

	% ----------------------------------------------------------------
	\subsection{A credit union's loan review (C1)}
	\label{prob:c1-5}

	A credit union wants to describe all 800 of its active microenterprise working-capital loans on June 30. In a random sample of 50 loans, the mean outstanding balance is 8.0 million COP and the sample standard deviation is 2.4 million COP. Ten of the sampled loans have an overdue installment on that date.

	\QuestionsTasks[sec:worked-solution-examples]{
		{1}/{sol:c1-5-a}/
		{Identify the population and sample. Give the symbols and contextual meanings for the population mean, standard deviation and proportion, and pair them with the corresponding sample statistics.},
		{1}/{sol:c1-5-b}/
		{The lending committee asks the analyst to include active microenterprise equipment loans in the June 30 review. The two loan categories together contain 1,000 loans. In a random sample of 50 from the combined list, the mean outstanding balance is 8.8 million COP and the standard deviation is 2.9 million COP. Eight sampled loans have an overdue installment on that date. Revisit your Task 1 answers. What changes, what stays the same, and which numerical differences can you calculate?}
	}

	\clearpage

	\Needspace{5\baselineskip}
	\phantomsection
	\label{sol:c1-5-a}
	\subsubsection*{C1 -- Match each quantity to its symbol}
	The population is \HighlightStroke{G11T2C1Population}{all 800 active} microenterprise working-capital loans on June 30. The sample is \HighlightStroke{G11T2C1Sample}{the 50} selected loans.

	\begin{center}
	\begin{tabularx}{\linewidth}{@{}*{2}{>{\raggedright\arraybackslash}X}@{}}
	\toprule
	Population parameter & Sample statistic \\
	\midrule
	$\HighlightStroke{G11T2C1PopulationMean}{\mu}$, mean outstanding balance across all 800 loans & $\HighlightStroke{G11T2C1SampleMean}{\bar{x}}=\HighlightStroke{G11T2C1SampleMean}{8.0}$ million COP, mean among the 50 loans \\
	$\HighlightStroke{G11T2C1PopulationSD}{\sigma}$, standard deviation of outstanding balances across all 800 loans & $\HighlightStroke{G11T2C1SampleSD}{s}=\HighlightStroke{G11T2C1SampleSD}{2.4}$ million COP, standard deviation among the 50 loans \\
	$\HighlightStroke{G11T2C1PopulationProportion}{p}$, fraction of all 800 loans with an overdue installment on June 30 & $\HighlightStroke{G11T2C1SampleProportion}{\hat p}=10/50=\HighlightStroke{G11T2C1SampleProportion}{0.20}$, fraction among the 50 loans \\
	\bottomrule
	\end{tabularx}
	\end{center}

	\secInicio[prob:c1-5][sec:worked-solution-examples]

	\Needspace{18\baselineskip}
	\phantomsection
	\label{sol:c1-5-b}
	\subsubsection*{C1 -- Compare what changes when a loan category is added}

	\begin{center}
	\begin{tabularx}{\linewidth}{@{}>{\raggedright\arraybackslash}p{.16\linewidth}*{2}{>{\raggedright\arraybackslash}p{.22\linewidth}}>{\raggedright\arraybackslash}X@{}}
	\toprule
	Quantity & Original scenario & Revised scenario & What the change tells us \\
	\midrule
	Population & \HighlightStroke{G11T2C1Population}{800 active} working-capital loans & \HighlightStroke{G11T2C1Population}{1,000 active} working-capital and equipment loans & The population gains $1{,}000-800=200$ loans. The review date stays the same. \\
	Sample & \HighlightStroke{G11T2C1Sample}{50 working-capital} loans & \HighlightStroke{G11T2C1Sample}{50 loans} from the combined list & The sample size stays at 50. The pool it is drawn from expands. \\
	Mean & $\HighlightStroke{G11T2C1PopulationMean}{\mu}$ for working-capital loans, $\HighlightStroke{G11T2C1SampleMean}{\bar{x}}=\HighlightStroke{G11T2C1SampleMean}{8.0}$ million COP & $\HighlightStroke{G11T2C1PopulationMean}{\mu}$ for both loan categories, $\HighlightStroke{G11T2C1SampleMean}{\bar{x}}=\HighlightStroke{G11T2C1SampleMean}{8.8}$ million COP & $8.8-8.0=+0.8$ million COP in $\HighlightStroke{G11T2C1SampleMean}{\bar{x}}$. Whether $\HighlightStroke{G11T2C1PopulationMean}{\mu}$ changed is unknown. \\
	Standard\newline deviation & $\HighlightStroke{G11T2C1PopulationSD}{\sigma}$ for working-capital loans, $\HighlightStroke{G11T2C1SampleSD}{s}=\HighlightStroke{G11T2C1SampleSD}{2.4}$ million COP & $\HighlightStroke{G11T2C1PopulationSD}{\sigma}$ for both loan categories, $\HighlightStroke{G11T2C1SampleSD}{s}=\HighlightStroke{G11T2C1SampleSD}{2.9}$ million COP & $2.9-2.4=+0.5$ million COP in $\HighlightStroke{G11T2C1SampleSD}{s}$. Whether $\HighlightStroke{G11T2C1PopulationSD}{\sigma}$ changed is unknown. \\
	Overdue share & $\HighlightStroke{G11T2C1PopulationProportion}{p}$ for working-capital loans, $\HighlightStroke{G11T2C1SampleProportion}{\hat p}=10/50=\HighlightStroke{G11T2C1SampleProportion}{0.20}$ & $\HighlightStroke{G11T2C1PopulationProportion}{p}$ for both loan categories, $\HighlightStroke{G11T2C1SampleProportion}{\hat p}=8/50=\HighlightStroke{G11T2C1SampleProportion}{0.16}$ & $0.16-0.20=-0.04$ in $\HighlightStroke{G11T2C1SampleProportion}{\hat p}$ ($-4$ percentage points). Whether $\HighlightStroke{G11T2C1PopulationProportion}{p}$ changed is unknown. \\
	\bottomrule
	\end{tabularx}
	\end{center}

	\secInicio[prob:c1-5][sec:worked-solution-examples]

	\clearpage
	\section{Advanced Problems}
	\label{sec:c1-advanced}

	\SectionProblemsTableOfContents{
		\SectionProblemLink{prob:c1-merchant-fees}{A retailer expands its fee review (C1)} \\
		\SectionProblemLink{prob:c1-settlement}{A payment processor checks its July report (C1)}
	}

	\subsection{A retailer expands its fee review (C1)}
	\label{prob:c1-merchant-fees}

	A retailer studies the processing fee per card payment in September. Its records contain 2,400 in-store card payments and 1,600 online card payments. A payment counts as high-fee if its processing fee is at least 2 thousand COP. The retailer draws three random samples as shown.

	\begin{center}
	\begin{tabularx}{\linewidth}{@{}l>{\raggedright\arraybackslash}Xrrrr@{}}
	\toprule
	\textbf{Draw} & \textbf{Records used for the draw} & \textbf{Size} & \textbf{Mean fee} & \textbf{Fee SD} & \textbf{High-fee count} \\
	\midrule
	A & In-store only & 60 & 1.4 & 0.6 & 12 \\
	B & In-store only, fresh draw & 60 & 1.6 & 0.7 & 15 \\
	C & In-store and online combined & 80 & 1.8 & 0.7 & 24 \\
	\bottomrule
	\end{tabularx}
	\end{center}
	Mean fees and fee standard deviations (SD) are in thousand COP.

	\QuestionsTasks[sec:c1-advanced]{
		{1}/{sol:c1-merchant-fees}/
		{Identify the sample and the target population for each draw. Label the three sample statistics and interpret their corresponding parameters. Calculate the changes from A to B and from B to C. Do the results establish that the combined population has a higher mean fee and high-fee share but the same fee spread as the in-store population?}
	}

	\subsubsection*{Complete worked solution}
	\label{sol:c1-merchant-fees}

	\textbf{Identify the groups and quantities.} A and B are separate \HighlightStroke{G11T2C1Sample}{samples of 60 in-store payments} from the same \HighlightStroke{G11T2C1Population}{2,400 September in-store payments}. C is a \HighlightStroke{G11T2C1Sample}{sample of 80 payments} drawn from \HighlightStroke{G11T2C1Population}{all 4,000 September card payments}, including both channels. For each target population, $\HighlightStroke{G11T2C1PopulationMean}{\mu}$ is the mean processing fee, $\HighlightStroke{G11T2C1PopulationSD}{\sigma}$ is its standard deviation, and $\HighlightStroke{G11T2C1PopulationProportion}{p}$ is the fraction of payments with fees of at least 2 thousand COP.

	\textbf{Compare the observed statistics.} The table in the question gives $\HighlightStroke{G11T2C1SampleMean}{\bar{x}}$ and $\HighlightStroke{G11T2C1SampleSD}{s}$. Dividing each high-fee count by its sample size gives
	\[
	\HighlightStroke{G11T2C1SampleProportion}{\hat{p}_A}=\frac{12}{60}=0.20,\qquad
	\HighlightStroke{G11T2C1SampleProportion}{\hat{p}_B}=\frac{15}{60}=0.25,\qquad
	\HighlightStroke{G11T2C1SampleProportion}{\hat{p}_C}=\frac{24}{80}=0.30.
	\]
	\begin{center}
	\begin{tabular}{@{}lcc@{}}
	\toprule
	\textbf{Statistic} & \textbf{B minus A} & \textbf{C minus B} \\
	\midrule
	$\HighlightStroke{G11T2C1SampleMean}{\bar{x}}$ (thousand COP) & $1.6-1.4=+0.2$ & $1.8-1.6=+0.2$ \\
	$\HighlightStroke{G11T2C1SampleSD}{s}$ (thousand COP) & $0.7-0.6=+0.1$ & $0.7-0.7=0$ \\
	$\HighlightStroke{G11T2C1SampleProportion}{\hat{p}}$ (percentage points) & $25-20=+5$ & $30-25=+5$ \\
	\bottomrule
	\end{tabular}
	\end{center}

	\textbf{Interpret the two comparisons.} A and B illustrate how statistics can differ while the population parameters stay fixed. From B to C, both the sample and the target population change. C has a higher observed mean and high-fee share than B, and the two samples have the same fee SD. The sample results leave the direction of the population differences in $\mu$ and $p$, and the equality of the two values of $\sigma$, undetermined.

	\secInicio[prob:c1-merchant-fees][sec:c1-advanced]

	\clearpage
	\subsection{A payment processor checks its July report (C1)}
	\label{prob:c1-settlement}

	A payment processor reviews card transactions made by its merchants in one city. Settlement time is the number of business days from payment until the merchant receives the funds. A transaction is late if settlement takes more than two business days. The complete June records give a mean settlement time of 1.8 business days and a late share of 16\%. In July, one random sample of 50 transactions has a mean of 1.6 business days and 5 late transactions. A second random sample of 50 from the same July list has a mean of 1.9 business days and 9 late transactions.

	\QuestionsTasks[sec:c1-advanced]{
		{1}/{sol:c1-settlement}/
		{An analyst writes, \lq\lq July's mean settlement time fell by 0.2 business days and the late share fell by 6 percentage points.'' Label June's and July's reported quantities and identify the populations and samples. Compare each July draw with June and then compare the two July draws. Rewrite the claim to reflect what is known and say what data would establish the exact month-to-month changes.}
	}

	\subsubsection*{Complete worked solution}
	\label{sol:c1-settlement}

	\textbf{Check the source of each number.} The two target populations are \HighlightStroke{G11T2C1Population}{all June transactions} and \HighlightStroke{G11T2C1Population}{all July transactions} for the processor's merchants in the city. June's complete records give $\HighlightStroke{G11T2C1PopulationMean}{\mu_{\text{June}}}=1.8$ business days and $\HighlightStroke{G11T2C1PopulationProportion}{p_{\text{June}}}=0.16$. Each July draw summarizes \HighlightStroke{G11T2C1Sample}{50 selected July transactions}. The means are $\HighlightStroke{G11T2C1SampleMean}{\bar{x}_1}=1.6$ and $\HighlightStroke{G11T2C1SampleMean}{\bar{x}_2}=1.9$ business days, and the late proportions are $\HighlightStroke{G11T2C1SampleProportion}{\hat{p}_1}=5/50=0.10$ and $\HighlightStroke{G11T2C1SampleProportion}{\hat{p}_2}=9/50=0.18$.

	\textbf{Calculate and label the differences.} Use the same definition of settlement time and lateness throughout.
	\begin{center}
	\begin{tabular}{@{}lcc@{}}
	\toprule
	\textbf{Comparison} & \textbf{Mean (business days)} & \textbf{Late share (percentage points)} \\
	\midrule
	July draw 1 minus June population & $1.6-1.8=-0.2$ & $10-16=-6$ \\
	July draw 2 minus June population & $1.9-1.8=+0.1$ & $18-16=+2$ \\
	July draw 2 minus July draw 1 & $1.9-1.6=+0.3$ & $18-10=+8$ \\
	\bottomrule
	\end{tabular}
	\end{center}

	\textbf{Revise the report.} \lq\lq Compared with the June population figures, the first July sample had a mean settlement time 0.2 business days lower and a late share 6 percentage points lower. The second July sample was 0.1 business days and 2 percentage points above the respective June figures.'' The analyst treated the first draw's statistics as July parameters. The draws differ even though they estimate the same $\HighlightStroke{G11T2C1PopulationMean}{\mu_{\text{July}}}$ and $\HighlightStroke{G11T2C1PopulationProportion}{p_{\text{July}}}$. To establish the exact differences between months, use the full July records or verified July population summaries, then subtract the corresponding June parameters.

	\secInicio[prob:c1-settlement][sec:c1-advanced]

\end{document}
